% ECONOMICS_PROBLEM: {"id": "F31-433", "title": "Dominant-currency pricing and exchange-rate pass-through", "jel_code": "F31", "source_id": "OP-0090"}
% !TeX program = xelatex
\documentclass[11pt,a4paper]{article}
\usepackage[margin=23mm]{geometry}
\usepackage{fontspec}
\setmainfont{Latin Modern Roman}
\usepackage{amsmath,amssymb,mathtools}
\usepackage{xcolor,enumitem,booktabs,longtable,array,xurl,hyperref}
\definecolor{muted}{HTML}{53636C}
\hypersetup{colorlinks=true,urlcolor=blue,linkcolor=blue}
\setlength{\parindent}{0pt}
\setlength{\parskip}{5pt}
\newcommand{\R}{\mathbb R}
\newcommand{\E}{\mathbb E}
\newcommand{\N}{\mathbb N}
\newcommand{\ind}{\mathbf 1}
\newcommand{\Var}{\operatorname{Var}}
\newcommand{\Cov}{\operatorname{Cov}}
\newcommand{\supp}{\operatorname{supp}}
\DeclareMathOperator*{\argmin}{arg\,min}
\DeclareMathOperator*{\argmax}{arg\,max}
\newcommand{\entrymeta}[3]{{\sffamily\small\color{muted}\textbf{\texttt{#1}}\quad|\quad #2\par #3\par}}
\newcommand{\Problemref}[1]{\texttt{#1}}
\newcommand{\JELref}[2]{\texttt{#2} (JEL \texttt{#1})}
\begin{document}
\section*{Dominant-currency pricing and exchange-rate pass-through}
\textbf{Problem ID:} \texttt{F31-433}\quad \textbf{JEL:} \texttt{F31}\par
% BEGIN ECONOMICS STATEMENT
\label{problem:OP-0090}
\label{atlas:prob:T10}

\noindent\textbf{Atlas ID:} T10; \textbf{original number:} 90; \textbf{field:} International trade. \textbf{Classification:} Editorial research formulation (theoretical/empirical); not a certified unsolved theorem.

\subsection*{Definitions and assumptions}
Let exporting firm $i$ choose an invoice currency $k_i$ from a finite set $K$, a reset price, inputs and financing arrangements. Its feasible choices depend on contracts, imported-input costs and currency liabilities. A dynamic model $\theta$ specifies demand, nominal rigidities, monetary policy, exchange rates and agents’ information. An equilibrium consists of optimal feasible firm and household decisions, a Nash equilibrium of currency choices conditional on the pricing environment, market clearing and a specified selection when currency complementarities create multiplicity. Let $p_{it}^d$ be the strictly positive destination-currency price and $e_t^{d,k}$ destination-currency units per unit of currency $k$.

\subsection*{Formal research question}
For an identified monetary-policy innovation $\varepsilon_t$ and horizon $h\le H$, estimate the equilibrium causal responses
\[
 B_{i,h}=\left.\frac{\partial}{\partial a}
 \mathbb E_\theta[\log p_{i,t+h}^{d}(\operatorname{do}(\varepsilon_t=a))]
 \right|_{a=0},
\]
along with exchange-rate, trade-volume and invoice-switching responses. Define pass-through as a ratio to the corresponding nonzero exchange-rate response, specifying the exchange-rate pair. Identify the determinants of equilibrium $k_i$, and characterize how alternative currency configurations alter the welfare and international transmission of monetary policy. If derivatives fail because currency switching is discrete, use finite policy changes.

\subsection*{Interpretation and scope}
A regression of prices on exchange rates does not automatically identify causal pass-through: demand, costs and policy respond jointly. Define whether prices are measured before or after a reset, and separate invoice-currency accounting from behavioral price adjustment. Changes in monetary policy may alter currency choice over longer horizons. An instrumental policy-surprise design requires an exclusion argument, especially when announcements reveal information about the economy. Currency denomination of liabilities and imported inputs can change optimal choices independently of market power. Welfare comparisons require household utility, financing and the monetary-policy rule rather than the price response alone. Fix a common initial state distribution for both policy interventions before taking their response difference.

\subsection*{Source audit and status}
[S1]–[S3] identify benchmark currency-pricing mechanisms and evidence. Atlas link [46], identified in [S4], studies macro-announcement price discovery in FX; it does not establish an invoicing equilibrium. None of these references certifies the entire editorial identification and policy question as unresolved in 2026.

\subsection*{References}

\begin{enumerate}[label={[S\arabic*]},leftmargin=*]

\item Gita Gopinath, Oleg Itskhoki, and Roberto Rigobon (2010). \emph{Currency Choice and Exchange Rate Pass-Through}. American Economic Review 100(1), 304–336. \url{https://doi.org/10.1257/aer.100.1.304}.
\textit{Locator and role:} Publisher or author abstract / bibliographic record; Invoice currency and pass-through of US import prices. Provides background or a benchmark result; does not establish that the editorial question remains unresolved in 2026.

\item Mary Amiti, Oleg Itskhoki, and Jozef Konings (2014). \emph{Importers, Exporters, and Exchange Rate Disconnect}. American Economic Review 104(7), 1942–1978. \url{https://doi.org/10.1257/aer.104.7.1942}.
\textit{Locator and role:} Publisher or author abstract / bibliographic record; Imported inputs, variable markups and exchange-rate pass-through. Provides background or a benchmark result; does not establish that the editorial question remains unresolved in 2026.

\item Gita Gopinath, Emine Boz, Camila Casas, Federico J. Díez, Pierre-Olivier Gourinchas, and Mikkel Plagborg-Møller (2020). \emph{Dominant Currency Paradigm}. American Economic Review 110(3), 677–719. \url{https://doi.org/10.1257/aer.20171201}.
\textit{Locator and role:} Publisher or author abstract / bibliographic record; Dominant-currency pricing, complementarities and imported inputs. Provides background or a benchmark result; does not establish that the editorial question remains unresolved in 2026.

\item Torben G. Andersen, Tim Bollerslev, Francis X. Diebold, and Clara Vega (2003). \emph{Micro Effects of Macro Announcements: Real-Time Price Discovery in Foreign Exchange}. American Economic Review 93(1), 38–62. \url{https://doi.org/10.1257/000282803321455151}.
\textit{Locator and role:} Publisher or author abstract / bibliographic record; High-frequency exchange-rate reactions to macroeconomic news. This source does not identify firms’ invoice-currency equilibrium or currency-switching responses.

\end{enumerate}

\subsection*{Common conventions: Source mathematical conventions}

Unless an entry states otherwise, agent and firm sets are finite. State and action spaces are measurable, strategies and outcome maps are measurable, probability laws are specified, and expectations exist for the targets under discussion. Infinite-horizon discount factors lie in $(0,1)$ and discounted objectives are finite. Norms, tolerances and complexity input sizes must be specified for computational comparisons. Constants or primitives that appear locally are inputs, not empirical facts asserted by this catalogue.

\textbf{Identification.} A statistical model $m\in\mathcal M$ implies an observable law $P_m$ and target $\theta(m)$. At law $P$, its identified set is
\[
\Theta(P;\mathcal M)=\{\theta(m):m\in\mathcal M,\ P_m=P\}.
\]
Point identification holds when this set is a singleton, or equivalently when $P_{m_1}=P_{m_2}$ implies $\theta(m_1)=\theta(m_2)$ for all compatible models. Sharp bounds mean bounds attained, or approached when the boundary is not attained, by compatible models. Writing this set is a definition, not a solution: the substantive task is to establish informative restrictions and derive or estimate the set. An unrestricted-confounding model may give only trivial bounds.

\textbf{Inference.} Identification, estimability and valid uncertainty quantification are separate questions. Fix $\alpha\in(0,1)$. A confidence procedure $C_n$ over a declared law class $\mathcal P$ has asymptotic uniform coverage at level $1-\alpha$ when
\[
\liminf_{n\to\infty}\inf_{P\in\mathcal P}\Pr_P\{\theta(P)\in C_n\}\geq1-\alpha.
\]
For set-identified targets the coverage event must be defined separately, for example coverage of the full identified set. If an entry uses identification-indexed models instead of uniquely indexed laws, the same coverage convention is applied over those models. Pointwise limits do not establish uniform validity, and asymptotic validity does not guarantee finite-sample precision. Sample sizes, dependence structures and weak-identification sequences are part of each application.

\textbf{Counterfactuals.} $Y_i(a)$ is a potential outcome under a specified intervention, including its timing, implementation and versions. It is not $Y_i$ conditional on voluntarily choosing $a$. A full intervention vector permits interference; shorthand scalar treatment notation does not silently impose no interference. Assignment, selection, attrition, anticipation and measurement must be specified. A claim about an instrument's local effect is not automatically a population policy effect.

\textbf{Economic equilibrium.} A counterfactual model includes best responses, constraints, feasibility, market clearing, state transitions and expectations consistent with the stated information. When several equilibria exist, either a declared selector is an input or the target is set-valued. Comparisons across policy regimes must use the stated selection convention. Reduced-form relationships are not presumed invariant after an equilibrium-changing intervention.

\textbf{Optimization and welfare.} Optimization is over the stated feasible set. If attainment is not established, $\sup$ or $\inf$ and $\varepsilon$-optimal choices replace an asserted maximizer or minimizer. For example, with value $V=\sup_{a\in\mathcal A}W(a)<\infty$, an $\varepsilon$-optimal set is $\{a\in\mathcal A:W(a)\geq V-\varepsilon\}$ for $\varepsilon>0$. Benefits, costs, risk and health or environmental outcomes can be aggregated only using declared units and conversion weights. Interpersonal utility weights, the treatment of fiscal transfers and foreign profits, discounting, and the behavioral welfare criterion are explicit inputs. A comparative-static derivative additionally requires existence and appropriate differentiability of the selected counterfactual.

\textbf{Sources.} Local labels [S1], [S2], and so forth restart in every question. Locators identify the relevant abstract, model, section or result at the level actually checked; a bibliographic or abstract-level verification is not represented as inspection of an entire proof. Working papers are distinguished from later publications and changing versions. Source summaries are paraphrased. All URL checks and editorial formulations refer to the audit date stated above.
% END ECONOMICS STATEMENT
\end{document}
